% \documentclass[paper=a4, fontsize=11pt]{scrartcl} % A4 paper and 11pt font size
\documentclass[11pt, a4paper]{book}
\usepackage[T1]{fontenc} % Use 8-bit encoding that has 256 glyphs
\usepackage[utf8]{inputenc}
\usepackage{fourier} % Use the Adobe Utopia font for the document - comment this line to return to the LaTeX default
\usepackage{listings} % para insertar código con formato similar al editor
\usepackage[spanish, es-tabla]{babel} % Selecciona el español para palabras introducidas automáticamente, p.ej. "septiembre" en la fecha y especifica que se use la palabra Tabla en vez de Cuadro
\usepackage{url} % ,href} %para incluir URLs e hipervínculos dentro del texto (aunque hay que instalar href)
\usepackage{graphics,graphicx, float} %para incluir imágenes y colocarlas
\usepackage[gen]{eurosym} %para incluir el símbolo del euro
\usepackage{cite} %para incluir citas del archivo <nombre>.bib
\usepackage{enumerate}
\usepackage{hyperref}
\usepackage{graphicx}
\usepackage{tabularx}
\usepackage{booktabs}
\usepackage{titlesec}

\setlength{\parskip}{2.2ex}
\titlespacing\section{0pt}{5.5ex plus 1ex minus .2ex}{1.5ex}\relax
\titlespacing\subsection{0pt}{5.5ex plus 1ex minus .2ex}{1.5ex}\relax


\lstset{frame=tb,
  language=Python,
  breaklines=true,
  showstringspaces=false,
  columns=flexible,
  numbers=none,
  commentstyle=\color{dkgreen},
  stringstyle=\color{mauve},
  tabsize=3
}

\usepackage[table,xcdraw]{xcolor}
\hypersetup{
	colorlinks=true,	% false: boxed links; true: colored links
	linkcolor=black,	% color of internal links
	urlcolor=cyan		% color of external links
}

\renewcommand{\familydefault}{\sfdefault}
\usepackage{fancyhdr} % Custom headers and footers
\pagestyle{fancyplain} % Makes all pages in the document conform to the custom headers and footers
\fancyhead[L]{} % Empty left header
\fancyhead[C]{} % Empty center header
\fancyhead[R]{Ángel Píñar Rivas} % My name
\fancyfoot[L]{} % Empty left footer
\fancyfoot[C]{} % Empty center footer
\fancyfoot[R]{\thepage} % Page numbering for right footer
%\renewcommand{\headrulewidth}{0pt} % Remove header underlines
\renewcommand{\footrulewidth}{0pt} % Remove footer underlines
\setlength{\headheight}{13.6pt} % Customize the height of the header

\usepackage{titlesec, blindtext, color}
\definecolor{gray75}{gray}{0.75}
\newcommand{\hsp}{\hspace{20pt}}
\titleformat{\chapter}[hang]{\Huge\bfseries}{\thechapter\hsp\textcolor{gray75}{|}\hsp}{0pt}{\Huge\bfseries}
\setcounter{secnumdepth}{4}
\usepackage[Lenny]{fncychap}


\begin{document}
	\begin{titlepage}
\newlength{\centeroffset}
\setlength{\centeroffset}{-0.5\oddsidemargin}
\addtolength{\centeroffset}{0.5\evensidemargin}
\thispagestyle{empty}

\noindent\hspace*{\centeroffset}\begin{minipage}{\textwidth}

\centering
\includegraphics[width=0.9\textwidth]{logos/logo_ugr.jpg}\\[1.4cm]

\textsc{ \Large TRABAJO FIN DE MÁSTER\\[0.2cm]}
\textsc{ GRADO EN INGENIERÍA INFORMÁTICA}\\[1cm]

{\Huge\bfseries My Many Reads \\}
\noindent\rule[-1ex]{\textwidth}{3pt}\\[3.5ex]
{\large\bfseries Plataforma social para lectores con sistema de recomendaciones }
\end{minipage}

\vspace{2.5cm}
\noindent\hspace*{\centeroffset}
\begin{minipage}{\textwidth}
\centering

\textbf{Autor}\\ {Ángel Píñar Rivas}\\[2.5ex]
\textbf{Director}\\ {Juan Julían Merelo Guervós}\\[2cm]
\includegraphics[width=0.3\textwidth]{logos/etsiit_logo.png}\\[0.1cm]
\textsc{Escuela Técnica Superior de Ingenierías Informática y de Telecomunicación}\\
\textsc{---}\\
Granada, Septiembre de 2022
\end{minipage}
\end{titlepage}

	\thispagestyle{empty}

\begin{center}
{\large\bfseries My Many Reads \\ Plataforma social para lectores con sistema de recomendaciones }\\
\end{center}
\begin{center}
Ángel Píñar Rivas\\
\end{center}

%\vspace{0.7cm}

\vspace{0.5cm}
\noindent{\textbf{Palabras clave}: \textit{software libre}, \textit{desarrollo ágil}, \textit{integración continua}, \textit{proceso de desarrollo}, \textit{control de calidad}, \textit{biblioteca}, \textit{gestión}, \textit{lecturas}, \textit{libros}, \textit{recomendaciones}, \textit{estadísticas}, \textit{popularidad}, \textit{tendencias}, \textit{usuarios}
\vspace{0.7cm}

\noindent{\textbf{Resumen}\\

El objetivo de este proyecto es definir un proceso de desarrollo claro y robusto, apoyándose en una mentalidad ágil y un conjunto de herramientas que garanticen la buena calidad del código y su documentación. Para poner estas prácticas en uso, se desarrollará un sistema de gestión de lecturas que ofrezca recomendaciones automáticas basadas en las lecturas de cada usuario además de la posibilidad de participar activamente en la plataforma creando recomendaciones personalizadas y otro contenido que pueda resultar de interés al resto de usuarios. Mediante la información y estadísticas de uso generadas en esta plataforma, se obtendrá el beneficio necesario para cubrir costes de mantenimiento.

\cleardoublepage

\begin{center}
	{\large\bfseries My Many Reads \\ Social platform for readers with recommendation system }\\
\end{center}
\begin{center}
	Ángel Píñar Rivas\\
\end{center}
\vspace{0.5cm}
\noindent{\textbf{Keywords}: \textit{open source}, \textit{agile development}, \textit{continuous integration}, \textit{development process}, \textit{quality assurance}, \textit{library}, \textit{management}, \textit{readings}, \textit{books}, \textit{recommendations}, \textit{statistics}, \textit{popularity}, \textit{trends}, \textit{users}
\vspace{0.7cm}

\noindent{\textbf{Abstract}\\

The main goal of this project is to define a clear and robust development process, relying on an agile mindset and a set of tools to ensure good quality of both code and documentation. To put these practices into use, a readings management system will be developed to offer automatic recommendations based on the readings of each user as well as the possibility to actively participate in the platform by creating personalized recommendations and other content that may be of interest to the rest of users. Through the information and usage statistics generated on this platform, the necessary profit will be obtained to cover maintenance costs.

\cleardoublepage

\thispagestyle{empty}

\noindent\rule[-1ex]{\textwidth}{2pt}\\[4.5ex]

D. \textbf{Juan Julián Merelo Guervós}, Profesor del Departamento de Arquitectura y Tecnología de Computadores de la Universidad de Granada.

\vspace{0.5cm}

\textbf{Informo:}

\vspace{0.5cm}

Que el presente trabajo, titulado \textit{\textbf{My Many Reads}},
ha sido realizado bajo mi supervisión por \textbf{Ángel Píñar Rivas}, y autorizo la defensa de dicho trabajo ante el tribunal
que corresponda.

\vspace{0.5cm}

Y para que conste, expiden y firman el presente informe en Granada a Septiembre de 2022.

\vspace{1cm}

\textbf{El director: }

\vspace{5cm}

\noindent \textbf{(Juan Julián Merelo Guervós)}

\chapter*{Agradecimientos}

A la resistencia, por el apoyo frente a las oleadas incesantes de presión y trabajo durante este máster, y la sensación de compañerismo en tiempos de pandemia. Sin vosotros, la sensación de soledad me habría llevado a abandonar este curso.

A Phoenix, por su empatía, su disposición a escuchar y la sensación de espalda cubierta, todo ello sin obligación de hacerlo.

A mi tutor, por aguantar mis cabezonerías con mi propia forma de trabajo, las revisiones frecuentes y las discusiones fructíferas que me han enseñado tanto información como puntos de vista.

A mis amigos y familia por escuchar y convertir mi agobio y preocupaciones en risas y diversión, cuidando de mi salud mental.

Y finalmente a mí mismo (aunque suene feo), por resistir la tentación de abandonar y seguir adelante hasta el final, aún con todo el agobio y estrés, y las motivaciones desaparecidas.




	\newpage
	\tableofcontents

	\newpage
	\listoffigures

	\listoftables 
	\newpage

	\chapter{Introducción}

A día de hoy, el afán por la lectura está aumentando entre no sólo los jóvenes (el sector que más preocupaba) sino adultos también. Una de las posibles causas de este hecho podría haber sido la pandemia, ya que con la cuarentena muchos han tenido que encontrar nuevos \textit{hobbies} a los que dedicarse desde su hogar para reemplazar actividades al aire libre.

Como con cualquier otro tipo de afición, no es sencillo iniciarse en el mundo de la lectura sin conocer un buen punto de partida, pues lo ideal es comenzar con algo fácilmente disfrutable, lo cual depende fuertemente de los gustos personales de la persona en concreto, por lo que en general se suele hacer una pequeña investigación sobre aquellos que podrían ser interesantes.

Este tipo de búsquedas no solo las realizan aquellas personas introduciéndose en la lectura, sino también los ya acostumbrados lectores que acaban con su lista de libros pendientes o prefieren descubrir libros nuevos. Al fin y al cabo, cualquier persona interesada en la lectura acaba buscando algo que leer por los medios de que disponga.

\section{Motivación}

Justo antes de la pandemia decidí explorar nuevos \textit{hobbies} que me separaran de la pantalla. Compré varios juegos de mesa, empecé a pintar miniaturas y decidí comprarme un libro con una premisa medianamente interesante. De estos tres, las miniaturas fueron las únicas que acabé descartando, porque aunque me gustan, son bastante caras, y como para los juegos de mesa dependes de organizarte con tus amigos, me enfoqué bastante en la lectura.

Me encontré con el problema de que no conocía a mucha gente que leyera, y además, que compartiera mis gustos, por lo que dependía de buscar en internet en base a lo poco que llevaba leído, y la verdad es que era una tarea algo tediosa. Se me ocurrió que podría existir una plataforma de \textit{reviews} de la que podría servirme, y me topé con \textit{Goodreads}, pero la verdad es que el sitio no terminaba de gustarme del todo. Conozco otras plataformas para series y películas que sí son bastante cómodas de usar para estas búsquedas, pero no para libros.

Finalmente me acabé acostumbrando bastante a mi método "manual" de buscar nuevos libros, introduje a un par de amigos con mis mismos gustos a la lectura para discutir con ellos y acabé disfrutando más del \textit{hobby} en general. Sin embargo, siempre pensé (y sigo pensando) que todo esto habría sido más fácil e incluso mejor, si hubiera una plataforma online en la que explorar libros, compartir opiniones y recomendaciones, y gestionar las lecturas, en cierto modo parecido a lo que ofrece \textit{Goodreads} pero eliminando todos los detalles, aspectos y funcionalidades que hacen que yo, a pesar de disfrutar mucho de la lectura y de llevarme bien con los ordenadores en general, no quiera usar esta plataforma que parece estar pensada para gente con mi misma afición, y que además es la más destacable de este ámbito con diferencia.

\section{Descripción del problema}

Existe una gran diversidad de géneros y gustos personales a tener en cuenta a la hora de realizar una buena elección de un libro. En el caso de no tener contactos que puedan aconsejar, habría que recurrir a otras fuentes de información, principalmente internet ya que es la más accesible, amplia y fácil de utilizar de las disponibles.

Se pueden encontrar ideas o sugerencias con simples búsquedas a través de Google, pero éstas suelen venir de revistas o blogs cuyos motivos de dichas recomendaciones no son transparentes, pudiendo ser, por ejemplo, totalmente personales del autor (una única persona) o haberse visto influenciados por un movimiento de patrocinio por parte de alguna editorial.

Esto no parece una solución ideal, pues un usuario promedio fácilmente puede desconfiar de estas recomendaciones, ya que por una parte se desconoce de dónde proceden y cuales son los motivos detrás de éstas, y por otra parte podrían no encajar con los gustos del usuario.

El propósito de este proyecto es acabar con este problema, permitiendo a todo aquel usuario con interés en la lectura servirse de unas recomendaciones totalmente personalizadas, apoyándose por opiniones que una multitud de usuarios similares aportan para cada libro. En otras palabras, elaborando las recomendaciones para los usuarios en base a sus lecturas y respaldando esas recomendaciones con opiniones de otros usuarios reales, por lo que se acaba creando una plataforma social en torno a la lectura en la que se refleja la opinión mayoritaria de sus usuarios, aportando mucha fiabilidad.

\section{Objetivos}
En el caso de \textit{My Many Reads}, el objetivo es construir una plataforma social en la que crear una comunidad de usuarios alrededor de la lectura. Partiendo de una base de datos de libros, se pretende reunir a una comunidad de usuarios con intereses en la lectura, ofreciéndoles una serie de servicios y funcionalidades como pueden ser:

\begin{itemize}
    \item Gestión de una biblioteca personal, en la que los usuarios puedan almacenar los libros leídos, pendientes y abandonados junto a una valoración.
    \item Sistema de recomendaciones automático construyendo los gustos de los usuarios a partir de sus bibliotecas personales.
    \item Sistema de búsqueda de libros por popularidad en la plataforma o sus características propias como pueden ser géneros, idioma o fecha de publicación.
    \item Foro de discusiones para cada libro en el que plasmar opiniones, teorías u otros pensamientos relacionados con dicho libro.
    \item Posibilidad de escribir recomendaciones personalizadas entre dos libros similares, que puedan ayudar a otros usuarios a tomar una decisión con su próxima lectura.
\end{itemize}

La meta global por tanto es la de construir una plataforma con características atractivas para atraer a un conjunto de usuarios. Sabiendo esto, el siguiente paso es establecer una lista de objetivos con el fin de planificar y diseñar en torno al cumplimiento de estos, enfocando todo el desarrollo en una dirección. En el caso de \textit{My Many Reads}, estos son los objetivos:

\begin{itemize}
    \item Cumplir con una funcionalidad mínima de consulta de libros y gestión listas o bibliotecas.
    \item Con el fin de que los usuarios puedran encontrar fácilmente buenas sugerencias para sus próximas lecturas, construir un sistema de recomendaciones que, a partir de los gustos del usuario, devuelva una lista de libros que podrían ser de interés para éste.
    \item Permitir a los usuarios participar activamente en creación de contenido de la plataforma mediante foros de discusiones, recomendaciones entre libros personalizadas y reviews.
    \item Servir como medio de comunicación para ofrecer noticias pertenecientes al mundo de la lectura.
    \item Recoger estadísticas de los usuarios de la plataforma para conocer información sobre, por ejemplo, las tendencias actuales.
\end{itemize}

Estos dos últimos objetivos tienen cierta dependencia a que el proyecto se encuentre desplegado y con una base mínima de usuarios. Para llegar a ese punto, es imprescindible cumplir los objetivos previos cuya intención es atraer usuarios a la plataforma.

	\chapter{Estado del arte}
\label{Estado del arte}

Con lo descrito hasta ahora, se pueden identificar fácilmente dos tipos de usuarios o clientes de \textit{My Many Reads}, aunque puede que en el futuro aparezcan más. Éstos son:
\begin{itemize}
    \item El \textbf{usuario lector}, que visita la plataforma como un punto de interés y apoyo para su \textit{hobby}.
    \item El \textbf{usuario editorial}, que consume el tráfico de información generado con el fin de introducir anuncios o hacer estudios de mercado.
\end{itemize}

Entre estos dos usuarios se puede generar la siguiente lista de problemas:
\begin{itemize}
    \item \textbf{Usuario lector}
    \begin{itemize}
        \item Necesita una herramienta para gestionar mis lecturas.
        \item Necesita recomendaciones de libros nuevos, que se adapten a mis gustos.
        \item Necesita explorar libremente los libros existentes por características como la popularidad, en caso de querer explorar nuevas opciones.
        \item Necesita conocer libros similares entre sí, junto a los motivos de la similitud, en caso de querer leer algo parecido a un libro que ya ha leido.
    \end{itemize}
    \item \textbf{Usuario editorial}
    \begin{itemize}
        \item Necesita saber cuales son las tendencias más populares, como autores o géneros, para conocer el estado del mercado actual.
        \item Necesita publicitar libros en lugares concurridos por la audiencia adecuada.
    \end{itemize}
\end{itemize}


Actualmente solo destacan otras dos plataformas similares, Goodreads y Librarything, siendo la primera la opción más conocida y popular en cuanto a este tipo de plataformas. A continuación se exploran las características de estas dos plataformas, listando aspectos positivos y negativos.

\section[Goodreads]{Goodreads\\ {\large \url{https://www.goodreads.com/}}}

Sin duda la opción más popular, esta plataforma ofrece una serie de opciones bastante llamativas y útiles:

\begin{itemize}
    \item Una biblioteca de libros a la que puedes añadir libros leídos, en progreso y en lista de espera, junto con tu puntuación, fechas de inicio y fin de lectura y \textit{review}.
    \item Permite añadir libros por tu cuenta, en caso de que éstos no se encuentren en la web de la página por ser ediciones antiguas, editoriales pequeñas o cualquier otro motivo.
    \item Un sistema de recomendaciones basado en tu biblioteca.
    \item Estadísticas de tus lecturas.
    \item Listas creadas por la comunidad que recogen libros con características comunes, como puede ser \textit{Books That Should Be Made Into Movies}.
    \item Un foro de discusiones, un apartado de citas famosas y secciones de AMA (\textit{Ask Me Anything}) dirigidas a autores.
    \item Una sección de noticias y artículos relacionados con los libros y el mundo de la lectura en general.
    \item Elementos de \textit{gamificación} de la lectura como preguntas diarias, un \textit{challenge} de lecturas anuales y cuestionarios acerca de libros.
    \item Permite a los usuarios escribir y compartir sus propias historias.
\end{itemize}

No por nada esta plataforma es la más popular, pero sin embargo no carece de aspectos mejorables:

\begin{itemize}
    \item La interfaz en sí es bastante mejorable, se hace difícil a un usuario promedio navegar y descubrir todo lo que ofrece sin tener que realizar un poco de esfuerzo.
    \item No permite explorar cómodamente la base de datos de libros, solo puedes ver las tendencias actuales (decididas por la plataforma mediante algún algoritmo desconocido) y explorar de un género concreto.
    \item Las listas mencionadas anteriormente las crea un único usuario, por lo que solo la opinión de una persona es la que dicta la relación entre todos los libros de la lista.
    \item El apartado de discusión de libros no ofrece mucha libertad, una vez lo encuentras parece más bien un conjunto de publicaciones filtrado según si son acerca de o relacionados con el libro concreto.
    \item Al visitar la página de un libro, si quieres consultar opiniones de la gente, lo que se te ofrece es su puntación media (no precisamente muy visible) y una sección de comentarios como la que te podrías encontrar en Amazon o cualquier tienda online. El foro de discusiones mencionado anteriormente es una fuente valiosa de \textit{feedback} que, aún con los problemas mencionados anteriormente, no parece ser accesible desde esta página.
    \item Información adicional como la existencia de adaptaciones en otros formatos (tales como películas o series) no aparece por ninguna parte.
    \item No se hace mención a diferentes ediciones de un libro ni a los idiomas en los que está disponible.
    \item El autor, otro elemento del cual un usuario podría usar para guiarse a la hora de buscar su siguiente lectura, solo dispone de un perfil muy similar al que podría tener cualquier usuario común, diferenciándose en una sección de AMA y de listado de sus libros. No parece haber un lugar en el que discutir sobre el estilo narrativo u otros elementos comunes a todos sus libros, que puedan llevar a un usuario a leer diferentes libros del autor que no compartan género entre sí.
\end{itemize}

\section[Librarything]{Librarything\\ {\large \url{https://www.librarything.com/}}}

La opción menos popular, no da indicios de la funcionalidad al completo que ofrece hasta que el usuario se registra. Además de ofrecer menos funcionalidad que Goodreads, parece que tiene algunos problemas adicionales. Comenzando por lo que ofrece esta plataforma:

\begin{itemize}
\item Una biblioteca de libros para cada usuario, junto con un sistema de recomendaciones automáticas en base a esta biblioteca.
\item Una sección de comunidades, similares a grupos de lectura, en las que compartir opiniones, embarcarse en retos o hacer concursos.
\item Un foro de discusiones, similar al de Goodreads, pero más accesible.
\item Búsqueda de autores.
\item Conexión de tu perfil con otras aplicaciones, como redes sociales o plataformas de intercambio de libros.
\item Muestra más metadatos sobre los libros que Goodreads, como idiomas en los que está disponibles y portadas que puedes encontrar.
\item Permite consultar estadísticas de los usuarios de la plataforma, lo cual resulta de utilidad a la hora de decidir si involucrarse con esta plataforma.
\end{itemize}

La funcionalidad ofrecida es claramente menor, aunque cabe decir que tiene lo suficiente como para poder atraer a un conjunto de usuarios bastante amplio. Pero de nuevo, no todo se ofrece de la mejor forma posible, existen aspectos susceptibles de mejora:

\begin{itemize}
\item De nuevo, la interfaz resulta claramente mejorable, lo que puede causar pérdida de usuarios.
\item Se usan diferentes bibliotecas según el estado de lectura de un libro, y solo permite al usuario añadir una puntuación y etiquetas a cada libro añadido a la biblioteca.
\item No existe un buscador libros o base de datos propia, lo que permite agregar dos veces el mismo libro si se busca en páginas o ediciones diferentes, por lo que aparecen inconsistencias.
\item Tiene un buscador de librerías o bibliotecas cercanas a tu zona, lo cual no está mal, pero existen servicios similares más populares y fiables, como el propio Google Maps.
\item Como en Goodreads, solo permite explorar libros por género o etiqueta, sin ningún tipo de filtro más allá de eso.
\item No existen listas de libros y las recomendaciones escritas por usuarios no tienen por qué incluir el motivo de la recomendación, lo cual no genera mucha confianza.
\end{itemize}

\section{Conclusiones}

Aunque también hay que mencionar que existen otras plataformas de este estilo pero con temática diferente, en cuanto a libros hay pocas opciones, y la que parece la mejor de ellas, Goodreads, tiene una amplitud de aspectos que se podrían mejorar, y por tanto introducirse en este entorno no parece una lucha imposible.

A continuación se muestran dos listas, una para enumerar todas las características que me han parecido muy buenas y que se encuentran presentes en alguna de estas plataformas, y otra para todos los aspectos negativos que considero que serían innecesarios o incluso perjudiciales para \textit{My Many Reads}.

Características más destacables:

\begin{itemize}
\item La posibilidad de crear una o varias colecciones (bibliotecas) de libros para recoger tus lecturas, con la posibilidad de añadir anotaciones para indicar el estado de la lectura u otras valoraciones.
\item Un sistema de recomendaciones automático en función de los libros incluidos en tu biblioteca.
\item Mantener una buena relación entre un libro y otros componentes relacionados con éste, como información adicional (idiomas disponibles, página del autor, fecha de publicación, ediciones...) o contenido generado por otros usuarios (recomendaciones personalizadas, \textit{reviews}, puntuación...)
\item Proveer de herramientas a los usuarios para que su feedback (opiniones, puntuaciones, recomendaciones, etc.) quede reflejado y sea de utilidad para el resto de usuarios y los algoritmos de la plataforma.
\end{itemize}

Aspectos negativos:

\begin{itemize}
\item La interfaz no refleja correctamente la importancia de algunos elementos que no deberían pasar muy desapercibidos, como aquellos que sirven para indicar la popularidad de un libro.
\item No se ofrece mucha libertad a la hora de explorar libros. Sería interesante una sección dedicada a buscar libros según géneros, popularidad, puntuación, etc. configurando algún tipo de filtro más complejo.
\item Al consultar los libros relacionados con uno en concreto, no especifican por qué están relacionados con éstos salvo por \textit{"A muchos usuarios también les gustó este otro libro"}. Más detalles del por qué están relacionados ayudarían a un usuario a tomar una decisión.
\item La abundancia de funcionalidades \textit{"Nice to have"} como los \textit{challenges}, \textit{quizzes}, etc. distraen al usuario de la funcionalidad principal.
\end{itemize}

	\chapter{Planificación}

En esta sección se recoge todo lo necesario para comprender cómo se planifica el desarrollo de este proyecto, describiendo las diferentes unidades de trabajo y sus características, enumerando las diferentes tareas planificadas y especificando la organización temporal del proyecto.

Así mismo, se explica la metodología seguida ya que de ésta depende, en parte, la planificación elaborada.

\section{Metodología utilizada}

Se pretende seguir una metodología de desarrollo ágil, elaborando tareas de pequeño tamaño y agrupándolas para cumplir en conjunto un objetivo de relativa importancia. Estas tareas se asignan a un \textit{sprint} de una duración concreta, permitiendo así tener claramente determinadas unidades de trabajo pequeñas y concretas, conformando funcionalidades u objetivos más complejos que se proponen cumplimentar en un intervalo de tiempo determinado.

Pero como en todo, pueden darse imprevistos durante el desarrollo que requieran de un tiempo adicional, por lo que esta planificación no es final en ningún momento, y siempre se debe contemplar la posibilidad de encontrarse con obstáculos no planificados que afecten al desarrollo correspondiente a un \textit{sprint}.

Para tener en cuenta apropiadamente estos imprevistos, como ya se ha mencionado anteriormente, no se crean y planifican tareas inmutables. A mayor avance en el tiempo más grande es la incertidumbre, por lo que se proponen una serie de objetivos poco detallados a cumplir en un largo periodo de tiempo, pero las tareas concretas para alcanzar estos objetivos solo se planifican en el momento en el que se vayan a realizar.

Como ya se comentará más adelante, el proyecto se está desarrollando en GitHub\cite{Github}, que ofrece herramientas que permiten elaborar una planificación, como \textit{issues} y \textit{milestones}, e incluso ofrece un tablero de KanBan y un \textit{backlog} de tareas en la sección de proyectos.

\subsection{Issues}

Las issues recogen el motivo u objetivo de un conjunto de cambios, tareas o acciones a realizar para contribuir al desarrollo o mantenimiento del proyecto. Por sí solas no ofrecen una gran versatilidad, pero gracias al sistema de etiquetas personalizables, a proyectos (para el tablero de KanBan y el backlog) y milestones empleadas en conjunto, se consigue un buen entorno de planificación.

La intención es que todos los cambios realizados en el proyecto (salvo quizás alguna excepción puntual) vayan asignados a un \textit{issue}, y que cada \textit{issue} esté asignado a un \textit{milestone} y/o haga referencia a otros issues en caso de considerar necesario subdividir en tareas aún más pequeñas.

Los tipos de \textit{issues} más destacables de mencionar son las historias de usuario (etiqueta \textit{user-story}). Éstas definen las historias de usuario, y por lo general están divididas en tareas, empleando otros issues que hagan referencia a la historia de usuario a la que van a contribuir.

\subsection{Milestones}

Cada \textit{milestone} representa un \textit{Producto Mínimamente Viable} o PMV. Un \textit{milestone} recoge todo el conjunto de tareas (\textit{issues}) necesarios de cumplimentar para alcanzar los requisitos propuestos del PMV propuesto. Del uso de estos \textit{milestones} como representaciones de PMVs se deduce fácilmente que se sigue una metodología que emplea iteraciones para diferentes entregas del proyecto en función de su desarrollo, es decir, cada nuevo \textit{milestone} se construye como un incremento del anterior, por lo que cada PMV debe suponer un incremento notable en la funcionalidad.

El hecho de que estos \textit{milestones} estén compuestos del conjunto de \textit{issues} necesarios para su terminación no quiere decir que absolutamente todos los \textit{issues} del proyecto estén asignados a un \textit{milestone}, puesto que como se ha explicado en la sección anterior, un \textit{issue} es una herramienta muy versátil, permitiendo instancias de éstos dedicadas a tareas cotidianas o de menor importancia que no encajen como parte de un \textit{milestone}, como por ejemplo la corrección de erratas.

Por lo general, un \textit{milestone} estará compuesto de historias de usuario o historias técnicas, pero eso no quiere decir que éstas (mayormente las historias de usuario) puedan formar parte de varios \textit{milestones}. Un ejemplo genérico de este caso sería el de la siguiente historia de usuario:

\begin{center}
    \textit{Como Usuario de My Many Reads,}\\
    \textit{quiero poder consultar recomendaciones personalizadas de libros,}\\
    \textit{para así ayudarme en la toma de decisión de mi próxima lectura.}
\end{center}

Esta historia se podría cumplimentar fácilmente con un algoritmo muy básico de recomendaciones en base a la lista de libros leídos por el usuario. Sin embargo, en el caso de \textit{My Many Reads}, esta historia necesita algo más que un algoritmo básico. Añadir complejidad al algoritmo para producir mejores resultados u ofrecer métodos alternativos de obtener recomendaciones también pueden ser parte de esta historia propuesta, pero no tendrían por qué serlo del mismo PMV, por lo que esta historia podría tener varios niveles de completitud, pudiendo pertenecer algunos de ellos a \textit{milestones} diferentes.

\section{Temporización}

Dado que se sigue una metodología ágil, no procede realizar una temporización completa del proyecto, si no que éste es un proceso que se realiza en varias ocasiones a lo largo del desarrollo, para periodos más cortos de tiempo, pero siempre con objetivos más difusos y abiertos al cambio en largos plazos de tiempo.

En primer lugar hay que categorizar los \textit{issues} de desarrollo del proyecto, para diferenciar entre varios tipos y asignar un tiempo a cada uno de éstos:

\begin{itemize}
    \item \textbf{Epopeyas.} Éstas representan un cambio notable en el proyecto. Son una agrupación de tareas, que se completa cuando la última de éstas se ha completado. Cada epopeya debe estar estimada a poder finalizarse en un periodo entre 5 y 8 días hábiles.
    \item \textbf{Tareas.} Unidades mínimas de trabajo. Por sí solas, pueden incluso no suponer ningún cambio en el proyecto. Se deben planificar para ser realizables en 1 a 3 días hábiles.
\end{itemize}

Con varias de estas epopeyas se definiría un \textit{milestone} que, como se indica en el apartado anterior, representa un PMV. Lo ideal sería que cada PMV se desarrollase en entre 4 y 6 semanas, lo que significa apróximadamente 4 epopeyas por PMV, pero es posible que en algunos casos se requierá más o menos tiempo o epopeyas para completar un \textit{milestone}.

Cabe mencionar que también se emplearán issues que sean excepciones a estas categorías. Éstas pueden consistir en pequeños arreglos independientes, sin epopeya asignada, u objetivos de mayor duración que no encajan con la definición de epopeya pero que tampoco llegan a representar un PMV, quizás pudiendo definirse como \textit{fase}. Existirán pocas excepciones que no encajen en las categorías definidas previamente, y es por eso por lo que se prefiere simplemente dejarlas como excepciones y no intentar categorizar todas ellas.

\section{Seguimiento del desarrollo}

En la página del repositorio en GitHub se encuentran todas las \textit{issues} y \textit{milestones} del proyecto, organizadas como ya se ha descrito anteriormente. Así mismo, en la sección de proyectos, se puede consultar un tablero \textit{KanBan} que contendrá tanto tareas como epopeyas de desarrollo y su estado actual, con el fin de apoyar la organización del desarrollo.

Dado que se están empleando los \textit{Pull Requests} para realizar los incrementos de progreso en el proyecto, éstos también son una fuente muy útil de información acerca del seguimiento y estado del desarrollo.

Adicionalmente, en la sección de \textit{Implementación} se enumerarán todos los PMV junto con sus epopeyas y las tareas que las forman, así como el tiempo dedicado a cada uno de estos elementos, que se irán actualizando conforme se desarrolle el proyecto.
	\chapter{Análisis del problema}

\section{Aprovechamiento de la plataforma}
Si bien es cierto que ya existe algunas plataformas similares (ver [estado del arte]()), se puede ver que no son perfectas, y de alguna manera se podrían introducir mejoras en algún sentido que podrían aumentar el número de usuarios, por lo que no sería descabellado pensar en hacerles competencia y tener algo de éxito.

El primer paso sería lograr una comunidad activa, para lo cual hay que ofrecer incentivos a los usuarios para emplear la plataforma, pues una red social sin usuarios no serviría para nada. Pero también hay que tener en cuenta que ofrecer un mero servicio a los usuarios sin obtener algún tipo de beneficio con el que cubrir los gastos de mantenimiento no lo más adecuado, por lo que habría que explotar este tráfico de usuarios e información de alguna manera:

\begin{itemize}
    \item Sirviendo como lugar publicitario para libros de todo tipo, ya que en la plataforma se concentrarían usuarios con interés en la lectura, y por tanto, posibles clientes.
    \item Así mismo, se podrían recoger datos sobre los gustos e intereses de los usuarios para que las propias editoriales sepan qué géneros y autores están más de moda, ayudándoles así a embarcarse en proyectos con mayor beneficio potencial.
\end{itemize} 

	\chapter{Implementación}

En este capítulo se describe el proceso de implementación junto con todos los elementos en los que éste se apoya, una descripción de todas las elecciones técnicas tomadas junto con su justificación.

En primer lugar se muestra el listado de las historias de usuario generadas junto con sus correspondientes tareas, que dan forma a las unidades de trabajo refinadas. Tras esto, se encuentran todos los dilemas tecnológicos considerados, junto con la solución seleccionada y  una lista de alternativas brevemente definidas con el fin de proveer una justificación fundada y válida de la elección tomada.

\subsection{Historias de usuario}

Partiendo de los requisitos previamente listados, refinándolos en historias de usuario, se puede dar paso a la elaboración de tareas que guíen la implementación del proyecto. Cabe aclarar que, como este desarrollo sigue una mentalidad ágil, las historias y tareas se han ido generando a corto plazo, por lo que la siguiente lista de historias de usuario y sus tareas se ha ido construyendo a lo largo del proceso de desarrollo.

Algunas historias, como ya se verá a continuación, aparecen repetidas. Esto se debe a que su desarrollo ha abarcado más de un PMV, pero en cada uno de estos se tratan diferentes tareas, por lo que solo se citaran las que correspondan al PMV descrito.

Así mismo cabe mencionar que el primer \textit{milestone} del proyecto se excluye de esta sección, pues está compuesto únicamente de documentación no directamente relacionada con las tareas de implementación.

\subsubsection{PMV-1. Estructura de datos básica y base de la lógica de negocio}

\begin{itemize}
    \item \href{https://github.com/Anglepi/My-Many-Reads/issues/7}{\textbf{HU-1}} Criterios de evaluación del jurado \\
    Como miembro del jurado, \\
    quiero recibir una documentación y presentación sobre este proyecto y su desarrollo, \\
    con el fin de poder evaluarlo en base a unos criterios específicos.
    \begin{itemize}
        \item \href{https://github.com/Anglepi/My-Many-Reads/issues/15}{\textbf{T-1.1}} Definir estructura básica de datos para los libros.
        \item \href{https://github.com/Anglepi/My-Many-Reads/issues/36}{\textbf{T-1.2}} Crear primera aproximación de BD de libros.
        \item \href{https://github.com/Anglepi/My-Many-Reads/issues/37}{\textbf{T-1.3}} Añadir requisitos para el segundo PMV.
    \end{itemize}
    \item \href{https://github.com/Anglepi/My-Many-Reads/issues/29}{\textbf{HU-2}} Solicitar información de libros \\
    Como Sergio, lector en MMR, \\
    quiero solicitar información sobre libros de una base de datos, \\
    con el fin de saber si me pueden interesar según sus características.
    \begin{itemize}
        \item \href{https://github.com/Anglepi/My-Many-Reads/issues/31}{\textbf{T-2.1}} Definir estructura básica de datos para los libros.
        \item \href{https://github.com/Anglepi/My-Many-Reads/issues/32}{\textbf{T-2.2}} Crear primera aproximación de BD de libros.
        \item \href{https://github.com/Anglepi/My-Many-Reads/issues/33}{\textbf{T-2.3}} Métodos para tratar las propiedades de los libros.
    \end{itemize}
    \item \href{https://github.com/Anglepi/My-Many-Reads/issues/30}{\textbf{HU-3}} Bibliotecas para gestionar lecturas
    Como Sergio, lector en MMR, \\
    quiero una biblioteca de libros de mi elección, \\
    con el fin de gestionar mis lecturas.
    \begin{itemize}
        \item \href{https://github.com/Anglepi/My-Many-Reads/issues/34}{\textbf{T-3.1}} Definir estructura básica de datos para las bibliotecas.
        \item \href{https://github.com/Anglepi/My-Many-Reads/issues/35}{\textbf{T-3.2}} Funcionalidad básica de gestión de bibliotecas.
    \end{itemize}
\end{itemize}

\subsection{Entorno de desarrollo}

A continuación se describen todos los elementos que conforman el entorno de desarrollo, justificando las elecciones tomadas enfrentándolas a las alternativas encontradas e indicando el por qué de la decisión en base a la naturaleza de este proyecto.

\subsubsection{Gestor de tareas}

\textbf{Elección final:} \href{https://www.gnu.org/software/make/}{\textit{make}}.

\textit{Make} es uno de los gestores de tareas más populares actualmente, muy probablemente debido a la sencillez de su uso e instalación y su versatilidad. El factor determinante para tomar esta decisión fue, principalmente, esta versatilidad, pues sabía que no solo lo usaría para gestionar actividades del desarrollo, sino para todo lo relacionado con la elaboración de este proyecto, incluida la documentación.

Esta elección fue de las primeras que hubo de tomarse, mucho antes de comenzar el desarrollo, pues como es lo habitual, el proyecto comenzó por una documentación e investigación que sirviera para justificar la viabilidad del mismo, sin ser muy relevantes los detalles técnicos del mismo.

Solo es necesario un fichero \href{https://github.com/Anglepi/My-Many-Reads/blob/main/Makefile}{\textit{Makefile}} en el que incluir las tareas a automatizar bajo sus nombres

Las alternativas a considerar se investigaron una vez se tomó la elección del lenguaje de programación a usar, ya que el resto de gestores de tareas genéricos con los que me he topado los descartaba, principalmente debido a la sintaxis que empleaban o a que requerían la instalación de complementos adicionales que, probablemente, no usaría para otra cosa, por lo que me pareció inconveniente.

\textbf{Alternativa principal:} \href{https://www.pyinvoke.org/}{\textit{invoke}}.

Es una muy buena herramienta para gestión de tareas en Python. Mediante un fichero \textit{tasks.py} se definen las tareas a realizar, y dado que la sintaxis de python es muy conocida y para nada compleja, no es difícil de construir. Además se pueden ejecutar órdenes en la línea de comandos muy fácilmente, por lo que también es muy versátil.

Muy probablemente habría sido la opción principal si se supiera desde el principio que se usaría Python para desarrollar el proyecto, dado que no tiene absolutamente nada que envidiarle a \textit{make}. Sin embargo éste se usaba desde el principio del proyecto, y dado que el cambio no aportaba ningún beneficio, se optó por descartar la opción.

\subsubsection{Editor de código}

\textbf{Elección final:} \href{https://code.visualstudio.com/}{\textit{Visual Studio Code}}.

Es el editor más potente y, desgraciadamente, eso se nota cuando se utiliza en máquinas de bajas prestaciones. Sin embargo, ofrece muchísimas opciones de configuración base, detecta una gran cantidad de lenguajes a los que provee de un formato visual de apoyo y es propiedad de \textit{Microsoft}, por lo que se espera un buen mantenimiento y actualizaciones frecuentes.

Destaca principalmente por ofrecer de forma nativa una terminal, que aunque por defecto usa la nativa del sistema en el que se encuentre instalado, es muy sencillo instalar y usar cualquier otra. Además existe una gama muy amplia de complementos gratuitos y desarrollados por la comunidad, muy fáciles de instalar y gestionar. 

El propio editor te ofrece sugerencias cuando detecta el entorno en el que estás trabajando, y ni si quiera es necesario recurrir al navegador web para consultar información de éstos o instalarlos, pues su gestor de complementos ya satisface todas estas necesidades. Algunos de estos complementos útiles para el proyecto son:
\begin{itemize}
    \item \href{https://marketplace.visualstudio.com/items?itemName=ms-python.python}{Python}, que ofrece características propias del lenguaje como \textit{linting}, \textit{debugging} y formateo de código.
    \item \href{https://marketplace.visualstudio.com/items?itemName=eamodio.gitlens}{GitLens}, que aumenta las funcionalidades relacionadas con Git en el editor.
\end{itemize}

\textbf{Alternativa principal:} \href{https://www.sublimetext.com/}{\textit{Sublime Text}}.

Es otro editor muy popular, y aunque no tenga un equipo tan grande como el de \textit{Visual Studio Code} detrás, recibe actualizaciones de manera frecuente con notables mejoras.

Una ventaja que tiene frente al anterior es que no consume tantos recursos, pero dado que tengo la suerte de trabajar en una máquina de buenas prestaciones, este aspecto realmente no me supone un beneficio.

Por otra parte, lo que me hizo descartar esta opción fue el proceso de búsqueda e instalación de complementos, que es más manual y tediosa, así como la ausencia de una terminal en el entorno de forma nativa.

Cabe destacar que existen otras alternativas, IDEs para Python, que también pueden considerarse, como \href{https://docs.python.org/3/library/idle.html}{\textit{IDLE}}, que viene integrado, o \href{https://www.jetbrains.com/es-es/pycharm/}{\textit{PyCharm}} de JetBrains, un equipo muy conocido por desarrollar muy buenos IDEs.

Un IDE sería ideal si solo se fuera a trabajar en código, pero dado que este proyecto también requiere de documentación y podría incluir una variedad de tecnologías a emplear, un simple editor de código con complementos es mucho más adecuado.

\subsubsection{Lenguaje de programación}

\textbf{Elección final:} \href{https://www.python.org/}{\textit{Python 3}}.

En un principio mi intención era tomar esta elección dándole bastante importancia a lo cómodo que me sentiría con un lenguaje en cuanto a gustos personales, ya que prefiero lenguajes tipados y más explícitos. Sin embargo, cuanto más avanzaba en las fases de diseño, más peso daba al sistema de recomendación. 

Es evidente que necesitaría algún tipo de API para servir los datos, pero por otra parte también necesitaría buenas herramientas de tratamiento y análisis de datos para poder construir un buen algoritmo de recomendaciones automático. Esto me llevó a pensar que probablemente el lenguaje que buscaba fuera bastante popular, pues habría tener buenas bibliotecas de desarrollo de APIs y de tratamiento de datos, las cuales suelen estar relacionadas con una buena comunidad de desarrolladores.

La decisión, como se indica arriba, fue Python. Este lenguaje he de admitir que no se adapta mucho a mis gustos personales en cuanto a sintaxis, pero lo que ofrece me será de gran utilidad para este desarrollo.

\textbf{Alternativa:} \href{https://go.dev/}{\textit{Go}}.

En cuanto a lenguajes existe mucha variedad, por la que la cantidad de alternativas que puede haber a Python, muchas de ellas totalmente viables a día de hoy, es considerable. Es posible que mi elección no haya sido la más acertada ya que me preocupa el rendimiento, pues en una asignatura del Master lo usé para construir un algoritmo genético y los tiempos de ejecución resultaron ser devastadores.

Es cierto que la velocidad del lenguaje \cite{languagespeed} es importante, pero también lo son las utilidades que ofrece y el nivel de mantenimiento y actualizaciones que prolonga su ciclo de vida.

Me he decantado por poner esta alternativa ya que es la que tenía en mente al comienzo del proyecto. Sintácticamente es un lenguaje que me encanta, además del control de punteros explícito y otras características, y ya lo había probado anteriormente para desarrollar una API REST y me gustó bastante. Además, tiene al equipo de Google detrás, por lo que inspira bastante confianza.

Demuestra ser bastante más rápido que Python en \href{https://benchmarksgame-team.pages.debian.net/benchmarksgame/fastest/go-python3.html}{bastantes situaciones} \cite{languagespeed}, pero por desgracia no encontré librerías prometedoras con utilidades para el tratamiento de datos que me permitieran, entre otras cosas, establecer correlaciones.

\subsubsection{\textit{Test framework.}}

\textbf{Elección final:} \href{https://docs.Pytest.org/}{\textit{Pytest}}

Existen muchos marcos de test diferentes, cada uno de ellos con fortalezas y debilidades además de otras características que aportan atractivo según la metodología y mentalidad con la que se enfoque el desarrollo.

En este caso se ha escogido \textit{Pytest} por la sencillez de su uso, ya que el único requisito que hay que aplicar en el código es prefijar el nombre de ficheros, clases y funciones con \textit{test}. Este prefijo es el que el \textit{framework}, a través de la orden \textit{Pytest}, buscará para determinar qué debe ejecutar, y el resultado de dichas ejecuciones determinará el éxito o fallo del test.

\textit{Pytest} por sí solo no ofrece utilidades para determinar la cobertura de los tests, es decir, las líneas de código y casos que ejecutan dichos tests, lo cual resulta de gran utilidad para determinar aquellas secciones o ramas de ejecución cuyo comportamiento no se está comprobando. Para ello se puede incluir \href{https://pypi.org/project/pytest-cov/}{un plugin} o instalar la herramienta \href{https://coverage.readthedocs.io/}{\textit{coverage}} por separado, ambas opciones siendo prácticamente similares.

\textbf{Alternativas:}

Una opción muy popular es \href{https://robotframework.org/}{Robot Framework}, mediante el cual el programador define procesos y les otorga nombres empleando \textit{keywords} en lenguaje natural, lo que le permite construir tests para funcionalidades más complejas ordenando estas \textit{keywords}, dando lugar a una lista de pasos en lenguaje natural que definen una funcionalidad completa.

Como salida ofrece, además del resultado en la línea de comandos, archivos en formato \textit{html} con la información en un formato gráfico más visual y organizado.

Todo esto, a primera vista, suena muy llamativo, pero implica una carga extra de trabajo y un aumento de la deuda técnica que, en este caso, no me resultaba nada de interés, por lo que la descarté.

\href{https://docs.python.org/3/library/unittest.html}{\textit{Python unittest}} es otra alternativa, muy similar a \textit{Pytest}, cuya diferencia principalmente es que requiere la importación de un módulo y los tests definidos deben ser estrictamente clases que hereden de una en el módulo importado, dentro de la cual se deben definir las funciones de \textit{testing}. Como la única diferencia notable es tener que escribir un poco más de código y no aporta más beneficios, \textit{unittest} quedó descartada.

\subsubsection{Herramientas ajenas a la programación.}

\textbf{Para escribir la documentación:} \href{https://www.latex-project.org/}{\textit{LaTeX}} junto a \href{https://www.markdownguide.org/getting-started/}{\textit{Markdown}}.

Aunque LaTeX es algo más tedioso de usar puesto que la configuración inicial es mucho más compleja que cualquier procesador de textos, es una opción muy potente a la hora de construir documentos bien organizados. Si tuviera que elegir una alternativa, probablemente sería \href{https://workspace.google.com/intl/es/products/docs/}{Documentos de Google}, no es la más potente en cuanto a características de formato comparado con \href{https://www.microsoft.com/es-es/microsoft-365/word}{Microsoft Word} o su contraparte libre \href{https://es.libreoffice.org/descubre/writer/}{Writer}, pero está en la nube de forma nativa, lo cual resulta muy cómodo para trabajar desde cualquier lugar.

Por otra parte está \textit{Markdown}. Se trata de un lenguaje de marcas para dar formato al texto muy popular, el cual se encuentra cada vez en más sitios web a disposición de sus usuarios. En el caso de este proyecto, utilizo \textit{Markdown} para toda la documentación que quiero que sea visible directamente en el \href{https://github.com/Anglepi/My-Many-Reads}{repositorio del proyecto}. Dado que no existen alternativas que estén tan extendidas, realmente no tiene competencia para su uso en sitios web.

\textbf{Revisión ortográfica:} \href{http://aspell.net/}{\textit{aspell}}.

Es un corrector ortográfico \textit{Open Source} que empleo para la corrección de mi documentación. El principal motivo de la elección de ésta es que, a parte de ser \textit{Open Source}, tiene soporte para ficheros de LaTeX. En el repositorio hay \href{https://github.com/Anglepi/My-Many-Reads/blob/main/scripts/spellcheck.sh}{un \textit{script}} que se encarga de su uso, apoyándose en un diccionario personalizado de \href{https://github.com/Anglepi/My-Many-Reads/blob/main/scripts/spellcheckDictionary.txt}{excepciones} ya que esta comprobación ortográfica es parte del proceso de integración, y si no supera esta comprobación no se permitirá la inclusión de los cambios a la rama principal.

La alternativa más potente es \href{https://www.antidote.info/en/}{\textit{Antidote} de \textit{Druide}}, que además ofrece comprobaciones gramaticales las cuales parecen ser bastante acertadas. Sin embargo, al ser una opción de pago, se ha descartado.

\subsection{Módulos importados}

\textbf{Biblioteca de aserciones:} \href{https://github.com/assertpy/assertpy}{\textit{assertpy}}

Una biblioteca con varias funciones de aserciones para encapsular todo el proceso de \textit{assert} junto con las comprobaciones pertinentes. Partiendo de la función \textit{assert\_that} y encadenando otras, permite construir aserciones en lenguaje casi natural.

La alternativa más prometedora que consideré fue \href{https://github.com/grappa-py/grappa}{\textit{Grappa}}, con la cual también construyes las aserciones formando frases en lenguaje natural encadenando funciones. No hay mucha diferencia con assertpy en cuanto a características ofrecidas, así que por gustos personales me decanté por la otra opción.
	\chapter{Conclusiones y trabajos futuros}

	
	\newpage
	\bibliography{biblio}
	\bibliographystyle{plain}
	
\end{document}

